\scanpage{30}
\uline{Corollaire 3} Soient \deleted{$f,g$} $X,Y$ deux schémas de type fini sur un corps $K$ ext. de type fini de $\mathbb Q$, \deleted{tels} \inserted{on suppose} que $Y$ se plonge dans \unsure{une|un} \deleted{\unsure{VA|V.A.}} extension d'une V.A. par un tore, $X$ réduit, \deleted{d'où} \inserted{$X,Y$, géom. $0$-connexes. Soit} $\overline K$ une clôture algébrique de $K$, d'où des extensions ``extérieures'' $\mathcal E_{X/K}$, \inserted{$\mathcal E_{Y/K}$,} de $\mathcal E_{\overline K/K}=\operatorname{Gal}(\overline K/K)$ par $\pi_1(\overline X)$, $\pi_1(\overline Y)$, et pour tt morphisme $f:X-Y$, un morphisme $\pi_1(f)$ de $\mathcal E_{X/K}$ \inserted{\deleted{\unsure{par|dans}} dans $\mathcal E_{Y/K}$.}

Soient $f,g:X=Y$ tels que $\pi_1(f)=\pi_1(g)$ extérieurement -- i.e. $\pi_1(f)$ et $\pi_1(g)$ \unsure{soient|sont} conjugués par un élément de $\pi_1(Y)$ -- alors $f=g$.

En fait, il suffit $=$ que les hom d'\uline{extensions} (``extérieures\unsure{''}, \inserted{donc d'extensions tt court!)} par $\pi_1(\overline X)_{\mathrm{ab}}$ et par $\pi_1(\overline Y)_{\mathrm{ab}}$ soient égaux, pour pouvoir conclure $f=g$. (C'est à dire, une fois que, \deleted{par} \inserted{grâce à} des hypothèses \uline{anabéliennes}, des hypothèses purement géométriques sur les actions de $\pi_1(\overline f)$, $\pi_1(\overline g)$ ont pu \unsure{de|être}
\margin{\deleted{$\unsure{f_{\mathrm{ab}}=g_{\mathrm{ab}}}$} \deleted{\unsure{induits:}} $\deleted{H_1(\unsure{\overline X|X},\mathbb Z_\ell)\ \unsure{\xrightarrow{f,g}}\ H_1(\unsure{\overline Y|Y|f},\mathbb Z_\ell)}$}
\margin{29}
